\documentclass[12pt,a4paper,twoside]{report}

\def\chapterstyletype{classic}

\newcommand{\AuthorName}{Marek Karbarz}
\newcommand{\SupervisorName}{dr Mariusz Startek}
\newcommand{\FieldName}{Inżynieria i analiza danych}
\newcommand{\StudyYear}{3}
\newcommand{\ThesisYear}{2026}
\newcommand{\ThesisTitle}{System webowy do scoringu i rankingowania instrumentów Giełdy Papierów Wartościowych na podstawie danych dziennych}
\newcommand{\ThesisTitleEnglish}{Web system for scoring and ranking Warsaw Stock Exchange instruments using end-of-day data}

\usepackage[left=3cm, right=2.5cm, top=2.5cm, bottom=2.5cm, headsep=1.2cm]{geometry}
\linespread{1.3}
\usepackage{polski}
\usepackage[utf8]{inputenc}
\usepackage[T1]{fontenc}
\usepackage{microtype}
\usepackage{pdflscape}
\usepackage{amsmath}
\usepackage{amsfonts}
\usepackage{amssymb}
\usepackage{amsthm}
\usepackage{graphicx}
\usepackage{subcaption}
\usepackage[all]{nowidow}
\emergencystretch=\maxdimen
\hyphenpenalty=10000
\hbadness=10000
\usepackage{xcolor}
\usepackage{textcomp}
\usepackage[official]{eurosym}
\usepackage[colorlinks=true,allcolors=black,plainpages=false]{hyperref}
\usepackage{indentfirst}
\usepackage{enumitem}
\setlist[enumerate]{noitemsep,nolistsep}
\setlist[itemize]{noitemsep,nolistsep}
\usepackage{tocloft}
\usepackage{tabularx}
\usepackage{multirow}
\usepackage{diagbox}
\usepackage{booktabs}
\usepackage{array}
\usepackage{caption}
\usepackage{lscape}
\usepackage[numbers]{natbib}
\usepackage{xurl}
\usepackage{titlesec}
\usepackage{ifthen}
\usepackage{ragged2e}

\renewcommand{\tabularxcolumn}[1]{>{\RaggedRight\arraybackslash}p{#1}}

\graphicspath{{assets/}}
\newcolumntype{Y}{>{\centering\arraybackslash}X}
\newcolumntype{j}{>{\raggedleft}p{200pt}}

\hypersetup{
    pdftitle={rankr - praca inżynierska},
    pdfauthor={\AuthorName}
}

\renewcommand*{\figurename}{Rys.}
\renewcommand*{\tablename}{Tab.}
\captionsetup{font=small, labelsep=period, justification=centerlast}

\renewcommand{\thesection}{\arabic{chapter}.\arabic{section}.}
\renewcommand{\thesubsection}{\arabic{chapter}.\arabic{section}.\arabic{subsection}.}
\setlength{\cftbeforechapskip}{0pt}
\renewcommand{\cftchapaftersnum}{.}

\ifthenelse{\equal{\chapterstyletype}{classic}}{
  \titleformat{\chapter}[display]
    {\normalfont\LARGE\bfseries\RaggedRight}
    {\chaptertitlename\ \thechapter}
    {0pt}
    {\LARGE}
}{%
  \titleformat{\chapter}[block]
    {\normalfont\LARGE\bfseries\RaggedRight}
    {}
    {0pt}
    {%
      \makebox[1em][l]{\thechapter.}%
      \hangindent=1em
      \hangafter=1
      \ignorespaces
    }
}

\titleformat{name=\chapter, numberless}[block]
  {\normalfont\LARGE\bfseries\RaggedRight}
  {}
  {0pt}
  {\hangindent=1em\hangafter=1}

\titleformat{\section}[hang]
  {\normalfont\Large\bfseries\RaggedRight}
  {\thesection}
  {0.3em}
  {}
\titleformat{\subsection}[hang]
  {\normalfont\large\bfseries\RaggedRight}
  {\thesubsection}
  {0.3em}
  {}
\titleformat{\subsubsection}[hang]
  {\normalfont\normalsize\bfseries}
  {}
  {0pt}
  {}

\titlespacing*{\chapter}{0pt}{0pt}{\baselineskip}

\begin{document}

\hypersetup{pageanchor=false}

\thispagestyle{empty}
\linespread{1.0}
\newgeometry{top=2.5cm, bottom=2.5cm, left=3cm, right=3cm}

\begin{titlepage}
    \def\dystans{1cm}
    \begin{center}
        \Large
        \begin{figure}[h]
            \centering
            \includegraphics[height=3.2cm]{prz_pl.png}%
            \hfill
            \includegraphics[height=3.2cm]{wmifs_pl.png}
        \end{figure}
        {\large\scshape Politechnika Rzeszowska}

        {\large im. Ignacego Łukasiewicza}

        {\large\scshape\MakeUppercase{Wydział Matematyki i Fizyki Stosowanej}}
    \end{center}
    \vspace{1.3cm}

    \begin{center}
        \begin{minipage}[c][5cm][c]{10cm}
            \centering
            {\Large \AuthorName} \par\vspace{0.4cm}
            {\Large\textbf{\textsc{\ThesisTitle}}} \par\vspace{0.4cm}
            Promotor: \SupervisorName
        \end{minipage}
    \end{center}

    \vspace{2cm}

    \begin{center}
        {\large \textsc{Praca inżynierska} \\
            \vspace{0.5cm}
            \large kierunek studiów: \FieldName \\
            \large rok studiów: \StudyYear}
    \end{center}
    \vspace{1.1cm}

    \vfill
    \begin{center}
        {\large Rzeszów \ThesisYear}
    \end{center}
\end{titlepage}
\restoregeometry
\linespread{1.3}

\newpage
\thispagestyle{empty}
~\vfill
\newpage

\hypersetup{pageanchor=true}
\renewcommand{\contentsname}{\LARGE Spis treści}
\setcounter{tocdepth}{1}
\tableofcontents
\cleardoublepage

\chapter*{Wstęp}
\addcontentsline{toc}{chapter}{Wstęp}

Porównanie spółek giełdowych wymaga połączenia informacji o ich finansach, cenach akcji i zachowaniu inwestorów. Dane te pochodzą z różnych źródeł i zmieniają się w różnym czasie. Celem pracy jest przygotowanie aplikacji webowej \textit{rankr}, która zbiera te informacje i porządkuje spółki Giełdy Papierów Wartościowych w Warszawie (GPW) według ustalonych reguł punktowych.

Projekt obejmuje koszyk 20 spółek zapisany w mapie WIG20 dołączonej do programu. Działają pobieranie i lokalny zapis danych, trzy rodziny punktacji oraz strona z rankingiem. Wynik ma służyć porównaniu względnego potencjału wzrostu w horyzoncie 3--12 miesięcy; przydatność tej interpretacji wymaga osobnej oceny. Zakres pracy nie obejmuje wykonywania transakcji ani integracji z brokerem.

% TODO(thesis): Complete the introduction after the evaluation and final scope are settled.

\cleardoublepage
\chapter{Analiza problemu i źródła danych}
\label{chap:data}

\section{Podstawowe pojęcia i cel porównania}

Akcja jest papierem wartościowym dającym udział we własności spółki. Inwestor może uzyskać dochód dzięki wzrostowi jej ceny oraz wypłacie części zysku, czyli dywidendzie \cite{stocks}. Cena jednej akcji sama w sobie nie mówi jednak, czy spółka jest atrakcyjna inwestycyjnie. Akcja kosztująca 20 zł nie musi być tańsza w stosunku do wyników przedsiębiorstwa niż akcja kosztująca 200 zł.

Zmianę ceny można przedstawić jako stopę zwrotu:
\begin{equation}
    r=\frac{P_1-P_0}{P_0},
\end{equation}
gdzie $P_0$ oznacza cenę początkową, a $P_1$ cenę końcową. Wzrost ze 100 do 110 zł daje $r=0{,}1$, czyli $10\%$. Ten prosty przykład pomija dywidendy i koszty transakcji. Pokazuje natomiast, dlaczego przy porównywaniu zmian cen potrzebna jest zmiana względna, a nie sama różnica w złotych.

W projekcie zastosowano scoring, czyli przypisywanie punktów według jawnych reguł. Ranking powstaje przez uporządkowanie spółek według sumy tych punktów. Uwzględniono trzy rodziny informacji: dane finansowe, historię cen oraz zmiany ujawnionych pozycji krótkich. Wyższy wynik oznacza korzystniejszą ocenę według przyjętych reguł, bez przeliczania punktów na przewidywany procent zysku.

\section{Zakres porównania i wymagania}

WIG20 jest indeksem opartym na portfelu akcji 20 największych i najbardziej płynnych spółek Głównego Rynku GPW. Indeks podsumowuje zmiany ich cen, a jego podstawowa wersja nie uwzględnia dochodu z dywidend \cite{gpwbenchmark}. W aplikacji nazwa WIG20 określa koszyk wejściowy. Sam poziom indeksu nie jest składnikiem punktacji spółek.

Lista używana przez program pochodzi z pliku \path{data/raw/wig20_symbols.csv} i jest dołączana podczas kompilacji. Program nie pobiera składu indeksu przy każdym uruchomieniu. Plik wiąże kod spółki z identyfikatorem ISIN, symbolem Yahoo Finance i klasyfikacją sektorową. ISIN służy do rozpoznania tego samego instrumentu w różnych źródłach, nawet gdy używają one innych nazw lub skrótów.

Dla aplikacji przyjęto następujące wymagania:
\begin{itemize}
    \item pobieranie danych bez kont i kluczy dostępowych do wykorzystywanych źródeł;
    \item lokalny zapis danych pozwalający ponowić obliczenia;
    \item rozróżnienie brakującej wartości i rzeczywistego zera;
    \item kontynuowanie pozostałych pobrań po błędzie pojedynczej spółki;
    \item prezentacja punktów trzech rodzin i wyniku końcowego w jednej tabeli.
\end{itemize}

Analiza korzysta z zakończonych sesji dziennych, czyli dni handlu na giełdzie. Nie wymaga obserwowania każdej transakcji w trakcie sesji. Taki zakres odpowiada przyjętemu horyzontowi porównania i ogranicza ilość pobieranych danych. Użytkownik uruchamia aplikację lokalnie, a wyniki przegląda w przeglądarce.

\section{Źródła danych}

Każda rodzina punktacji ma oddzielne źródło danych. Tabela~\ref{tab:sources} przedstawia źródła wykorzystywane przez działający program. Mapa spółek określa, które instrumenty należy pobrać, natomiast wartości finansowe i rynkowe pochodzą z serwisów zewnętrznych. Dane o pozycjach krótkich udostępnia Komisja Nadzoru Finansowego (KNF).

\begin{table}[htbp]
\centering
\caption{Źródła danych wykorzystywane przez aplikację}
\label{tab:sources}
\small
\begin{tabularx}{\textwidth}{p{3cm}XX}
\toprule
Źródło & Zakres danych & Zastosowanie \\
\midrule
Mapa WIG20 w projekcie & Kody, ISIN, symbole Yahoo i sektory & Lista spółek i powiązanie danych \\
GPW / Notoria & Bieżąca tabela danych finansowych & Wskaźniki i punkty fundamentalne \\
Yahoo Finance & Dzienne ceny i wolumen & Sezonowość oraz trend \\
Publiczny rejestr KNF & Ujawnione pozycje krótkie netto & Punkty zmiany zaangażowania krótkiego \\
\bottomrule
\end{tabularx}
\end{table}

W mapie zapisano również sektor według klasyfikacji GICS. Jest to system grupowania przedsiębiorstw według rodzaju ich działalności \cite{gics}. Program wykorzystuje sektor do wyboru zestawu wskaźników fundamentalnych i progów punktowych. Nie wyznacza na tej podstawie osobnych rankingów branżowych.

\section{Historia cen}

Dla każdego dnia notowań zapisywane są cena otwarcia, najwyższa i najniższa cena, cena zamknięcia oraz wolumen, czyli liczba akcji objętych obrotem. Zestaw ten określa się skrótem OHLCV. Jedną taką dzienną obserwację nazywa się świecą. Historia świec pozwala porównać bieżącą cenę ze średnią z wcześniejszych sesji oraz obliczyć zmiany cen w poszczególnych miesiącach.

Yahoo Finance dostarcza również skorygowaną cenę zamknięcia, uwzględniającą odpowiednie podziały akcji i wypłaty dywidend \cite{yahoo}. Podział akcji zwiększa ich liczbę przy odpowiednim zmniejszeniu ceny jednostkowej. Zachowanie osobnego pola ceny skorygowanej pozwala korzystać z korekt dostawcy bez samodzielnego odtwarzania tych zdarzeń. Program używa tej ceny do sezonowości, a zwykłego pola zamknięcia do oceny trendu.

Przy pierwszym pobraniu program żąda całej dostępnej historii dziennej spółki. Kolejne uruchomienia pobierają zakres po ostatniej zapisanej świecy i dopisują zakończone sesje. Parser, czyli część programu odczytująca odpowiedź dostawcy, sprawdza zgodność instrumentu oraz poprawność cen. Nieprawidłowe lub nieukończone świece są pomijane z zapisaniem przyczyny, bez uzupełniania ich poprzednią ceną.

\section{Dane finansowe}

Dane finansowe opisują różne strony działalności przedsiębiorstwa. Rachunek wyników przedstawia przychody, koszty i wynik za określony okres. Bilans pokazuje stan aktywów, czyli zasobów przedsiębiorstwa, oraz źródła ich finansowania na dany dzień. Rachunek przepływów pieniężnych przedstawia wpływy i wydatki, dlatego uzupełnia informację o samym zysku \cite{statements}.

Importer odczytuje bieżącą tabelę GPW / Notoria wraz z okresem raportu, walutą i mnożnikiem jednostki. Mnożnik informuje na przykład, że kwoty podano w tysiącach złotych. Program zachowuje też oznaczenie, czy raport dotyczy pojedynczego podmiotu, czy grupy kapitałowej. Jeden zapisany stan tabeli jest dalej nazywany obserwacją lub snapshotem.

Punktacja wykorzystuje relacje między kwotami z tej samej obserwacji. Przykładem jest marża netto, czyli stosunek wyniku netto do przychodów. Przy zysku 10 mln zł i przychodach 100 mln zł taka relacja wynosi $10\%$. W użytym modelu wynik netto i kapitał własny dotyczą udziałowców jednostki dominującej, czyli podmiotu kontrolującego grupę.

Brak pozycji źródłowej oznacza brak odpowiedniego wskaźnika. Ma to znaczenie zwłaszcza dla banków, których dane różnią się od typowej tabeli przedsiębiorstwa niefinansowego. Nazwa okresu, na przykład „I--II kw. 2026”, pozostaje etykietą raportu. Nie jest traktowana jako data jego publikacji.

\section{Ujawnione pozycje krótkie}

Pozycja krótka pozwala co do zasady osiągać korzyść przy spadku ceny instrumentu. W pracy wykorzystano publiczny rejestr KNF, w którym dla spółki podawane są między innymi posiadacz pozycji, jej wielkość i data obliczenia. Ujawnienie publiczne dotyczy pozycji od $0{,}5\%$ wyemitowanego kapitału podstawowego spółki na posiadacza \cite{knfshorts}. Dane są powiązane ze spółkami przez ISIN.

Program pobiera strony rejestru wspólnie dla potrzebnych spółek, a następnie wybiera pasujące wpisy lokalnie. Dla jednego posiadacza zachowuje najnowszą pozycję, aby nie dodawać kolejnych zgłoszeń tej samej pozycji. Porównanie sum z dwóch obserwacji jest w projekcie uproszczonym sygnałem sentymentu, czyli nastawienia inwestorów. Zakres interpretacji tego sygnału omówiono w rozdziale~\ref{chap:tests}.

\cleardoublepage
\chapter{Architektura systemu}
\label{chap:architecture}

\section{Organizacja aplikacji}

Aplikacja powstała w języku Rust i składa się z jednego pakietu o nazwie \texttt{rankr}, przechowywanego w katalogu \texttt{importer/}. Pakiet zawiera bibliotekę z funkcjami pobierania i obliczeń oraz program wykonywalny, który uruchamia cały proces. Rust pozwala umieścić bibliotekę i program w jednym pakiecie \cite{rust}. Dzięki temu ten sam kod obliczeniowy może być używany przez stronę rankingu i testy.

Importer i serwer działają w jednym procesie. Alternatywą byłoby uruchamianie osobnego programu zbierającego dane i osobnej usługi prezentującej wyniki. Wspólny program upraszcza lokalne uruchomienie oraz przekazywanie danych między jego częściami. Konsekwencją jest oczekiwanie na zakończenie odświeżania przed udostępnieniem strony.

Serwer to część aplikacji odpowiadająca na żądania przeglądarki. W projekcie tę rolę pełni Axum, biblioteka obsługująca żądania HTTP i kierująca je do odpowiednich funkcji \cite{axum}. HTTP jest protokołem wymiany żądań i odpowiedzi między przeglądarką a serwerem. Przeglądarka otrzymuje gotowy dokument HTML opisujący zawartość tabeli.

\section{Przepływ danych}

Przebieg działania przedstawiono na rysunku~\ref{fig:pipeline}. Najpierw program sprawdza lokalne pliki i pobiera potrzebne aktualizacje. Dopiero po tej części uruchamia serwer pod lokalnym adresem. Wejście na stronę uruchamia obliczenie rankingu na podstawie zapisanych danych.

\begin{figure}[htbp]
\centering
\begin{tabular}{c}
\fbox{Lista spółek i sprawdzenie lokalnych danych} \\
$\downarrow$ \\
\fbox{Pobranie, odczyt i zapis aktualizacji} \\
$\downarrow$ \\
\fbox{Uruchomienie lokalnego serwera} \\
$\downarrow$ \\
\fbox{Żądanie przeglądarki: obliczenie rankingu} \\
$\downarrow$ \\
\fbox{Odpowiedź HTML z tabelą}
\end{tabular}
\caption{Kolejność działania aplikacji rankr}
\label{fig:pipeline}
\end{figure}

Pobieranie zależy od dostępności serwisów zewnętrznych, natomiast obliczenia korzystają już z plików na dysku. Odświeżenie strony ponownie liczy wynik, ale nie rozpoczyna kolejnego pobierania. Ponowną aktualizację danych wykonuje następne uruchomienie programu. Ranking dotyczy bieżącej daty UTC, czyli daty według uniwersalnego czasu stosowanego przez aplikację.

Do obsługi operacji sieciowych wykorzystano Tokio. Jest to środowisko wykonujące zadania asynchroniczne, które mogą oddać sterowanie podczas oczekiwania, na przykład na odpowiedź serwera \cite{tokio}. Odczyt plików i obliczenia rankingu są zlecane osobnemu wątkowi, czyli oddzielnemu ciągowi wykonywania zadań wewnątrz programu. Pozwala to oddzielić tę pracę od obsługi żądania sieciowego. Same historie cen są pobierane kolejno dla poszczególnych spółek.

\section{Lokalny zapis danych}

Dane są przechowywane w katalogu \path{data/collected}, liczonym względem miejsca uruchomienia programu. Zastosowano JSON, czyli tekstowy format zapisujący między innymi nazwane pola, liczby i listy. Biblioteka Serde JSON pozwala odczytywać takie dokumenty do struktur Rust i zapisywać struktury ponownie do plików \cite{serdejson}. Ułatwia to porównanie odpowiedzi dostawcy z wynikiem jej przetwarzania.

Przykładowy układ katalogów ma postać:
\begin{verbatim}
data/collected/
  YYYY-MM-DD_KGHM_fundamentals/
    raw.json
    fundamentals.json
  YYYY-MM-DD_KGHM_history/
    raw.json
    history.json
  YYYY-MM-DD_KGHM_shorts/
    raw.json
    shorts.json
\end{verbatim}

Plik \texttt{raw.json} przechowuje treść odpowiedzi oraz informacje o jej źródle i pobraniu. Fundamenty i dane KNF są archiwizowane przed interpretacją odpowiedzi; błąd odczytu pozostawia również opis problemu. Dla cen zachowywane są odpowiedzi zawierające nowe świece oraz odpowiedzi, których nie udało się poprawnie odczytać. Pusta aktualizacja cen nie tworzy nowych plików.

Alternatywą dla osobnych plików jest lokalna baza SQLite, która przechowuje tabele i indeksy w jednym pliku i nie wymaga osobnego serwera \cite{sqlite}. JSON pozwala jednak bezpośrednio przeglądać zapisane obserwacje i korzystać z istniejących struktur programu. Kosztem tego rozwiązania jest przeszukiwanie katalogów przy każdym obliczeniu rankingu. Aplikacja nie wymaga obecnie połączenia z bazą danych.

\section{Zasady aktualizacji i obsługa błędów}

Fundamenty i pozycje krótkie są sprawdzane osobno dla każdej spółki. Jeżeli istnieje odpowiedni plik w katalogu oznaczonym bieżącym dniem UTC, ponowne pobranie zostaje pominięte. W innym przypadku powstaje kolejna obserwacja, nawet jeżeli treść raportu nie zmieniła się od poprzedniego dnia. Historia cen jest natomiast sprawdzana przy każdym uruchomieniu i rozszerzana o późniejsze zakończone sesje.

Kolejne obserwacje fundamentów i KNF pozostają oddzielnymi plikami. Wybrany plik historii cen jest uzupełniany w miejscu, bez ponownego zapisywania całej serii w kolejnym katalogu. Powtórzone daty nie są dopisywane, a zapis kompletnego pliku zastępuje poprzednią wersję atomowo. Oznacza to, że nieudana próba zapisu nie powinna pozostawić częściowo zapisanej historii.

Błąd jednej spółki jest wypisywany w terminalu, po czym program przechodzi do pozostałych pobrań. Wcześniejsze poprawne dane pozostają dostępne. Serwer uruchamia się po zakończeniu całej próby odświeżania, również gdy część pobrań się nie powiodła. Stanowi to kompromis między dostępnością strony a aktualnością danych; reguły braków i świeżości wyniku opisano w rozdziale~\ref{chap:scoring}.

\section{Podział odpowiedzialności i prezentacja}

Kod rozdziela pobieranie, interpretację odpowiedzi, zapis, obliczenia i prezentację. Parser otrzymuje dokument, a moduł punktacji otrzymuje odczytane wartości. Takie funkcje można sprawdzać bez połączenia z serwisem zewnętrznym. Tabela~\ref{tab:modules} przedstawia główne grupy modułów, bez wyliczania wszystkich plików.

\begin{table}[htbp]
\centering
\caption{Główne części programu}
\label{tab:modules}
\small
\begin{tabularx}{\textwidth}{p{3cm}XX}
\toprule
Część & Odpowiedzialność & Przykładowy moduł \\
\midrule
Pobieranie & Sprawdzenie aktualności i żądania do źródeł & \texttt{data\_refresh} \\
Odczyt i zapis & Interpretacja odpowiedzi oraz pliki lokalne & \texttt{parser}, \texttt{storage} \\
Punktacja & Wyniki rodzin i połączenie ich w ranking & \texttt{ranking} \\
Prezentacja & Odpowiedź dla przeglądarki i tabela HTML & \texttt{page} \\
\bottomrule
\end{tabularx}
\end{table}

Strona pokazuje kod spółki, sektor, trzy wyniki rodzin i ich sumę. Pusta komórka oznacza brak wyniku, a liczba zero pozostaje widoczna jako \texttt{0}. Serwer sam przygotowuje HTML i prosty arkusz stylów określający wygląd tabeli. Przeglądarka odpowiada za wyświetlenie tabeli, a obliczenia odbywają się na serwerze.

% TODO(thesis): Revisit the architecture after the final interface and persistence scope are settled.

\cleardoublepage
\chapter{Model danych}
\label{chap:model}

\section{Instrument i obserwacje}

Spółkę identyfikują kod i ISIN. Każde pobranie ma zapisane źródło oraz czas \texttt{fetched\_at}; czas pobrania pozostaje odrębny od daty sesji i okresu raportu. Podstawowe dokumenty zestawiono w tabeli~\ref{tab:documents}.

\begin{table}[htbp]
\centering
\caption{Podstawowe pliki danych}
\label{tab:documents}
\small
\begin{tabularx}{\textwidth}{p{4cm}X}
\toprule
Plik & Zawartość \\
\midrule
\texttt{raw.json} & Odpowiedź źródłowa, adres, status HTTP i czas pobrania \\
\texttt{fundamentals.json} & Dane finansowe, okres, waluta i jednostka \\
\texttt{history.json} & Dzienne świece i przyczyny pominięcia niepoprawnych obserwacji \\
\texttt{shorts.json} & ISIN, czas pobrania i lista ujawnionych pozycji \\
\bottomrule
\end{tabularx}
\end{table}

\section{Liczby i brakujące wartości}

Kwoty i ceny są przechowywane w typie \texttt{Decimal}, który zachowuje dziesiętny zapis liczb bez pośredniej konwersji do binarnego typu zmiennoprzecinkowego \cite{decimal}. Brak wartości jest zapisywany jako \texttt{null}, a przy prezentacji wyniku jako pusta komórka.

% TODO(thesis): Add a compact data relationship diagram and examples of the final document fields.

\cleardoublepage
\chapter{Metodyka scoringu}
\label{chap:scoring}

\section{Punkty fundamentalne}

Program oblicza pięć relacji z jednego raportu i przyznaje punkty według tabeli~\ref{tab:thresholds}, w której $x$ oznacza wartość wskaźnika. Progi są przyjętą konwencją projektu. Wartości między progami, wraz z granicami, dają zero, a wynik rodziny jest sumą dostępnych punktów.

\begin{table}[htbp]
\centering
\caption{Progi punktacji poza sektorem Financials}
\label{tab:thresholds}
\small
\begin{tabularx}{\textwidth}{Xp{3cm}p{3cm}}
\toprule
Wskaźnik & $+1$ & $-1$ \\
\midrule
Marża netto: wynik netto / przychody & $>0{,}02$ & $<-0{,}02$ \\
Marża operacyjna: wynik operacyjny / przychody & $>0{,}02$ & $<-0{,}02$ \\
Udział kapitału: kapitał własny / aktywa & $>0{,}50$ & $<0{,}30$ \\
Zobowiązania / kapitał własny & $0\leq x<0{,}50$ & $x<0$ lub $x>1{,}50$ \\
Przepływy operacyjne / wynik netto & $>1$ & $<0$ \\
\bottomrule
\end{tabularx}
\end{table}

Relacja zobowiązań do kapitału wymaga obu części zobowiązań: długo- i krótkoterminowej. Relacja przepływów do wyniku jest liczona tylko przy dodatnim wyniku netto. Brak potrzebnej kwoty lub zerowy mianownik daje \texttt{null}; brak wszystkich punktów oznacza brak wyniku rodziny.

Dla sektora finansowego (Financials) punktowany jest tylko udział kapitału: powyżej $0{,}10$ otrzymuje $+1$, poniżej $0{,}08$ otrzymuje $-1$, a pomiędzy nimi zero. Reguła obejmuje cały sektor z mapy, również PZU i KRUK. Pozostałe cztery wskaźniki są w tym szablonie pomijane.

% TODO(thesis): Expand the interpretation of the ratios and justify the adopted threshold convention.

\section{Punkty techniczne}

Wynik techniczny jest sumą punktów sezonowości i trendu:
\begin{equation}
    T=S+G.
\end{equation}
Symbol $S$ oznacza punkty sezonowości, a $G$ punkty trendu; każdy składnik przyjmuje $-1$, $0$ lub $+1$. W obliczeniach uczestniczą tylko świece wcześniejsze niż data rankingu.

Sezonowość wykorzystuje średnią zmian cen w tym samym miesiącu ze wszystkich dostępnych wcześniejszych lat. Zwrot miesiąca jest liczony między ostatnimi dostępnymi sesjami tego miesiąca i poprzedniego, z ceny skorygowanej, a przy jej braku ze zwykłego zamknięcia. Średnia powyżej $2\%$ daje $+1$, poniżej $-2\%$ daje $-1$, a pozostałe wartości zero. Wystarcza jeden poprawny historyczny miesiąc; brak takich obserwacji oznacza brak wyniku.

Trend porównuje ostatnie zamknięcie ze SMA200, czyli średnią arytmetyczną 200 ostatnich dziennych zamknięć. Odchylenie mniejsze niż $3\%$ daje zero, a od $3\%$ cena powyżej średniej daje $+1$ i poniżej $-1$. Przy mniej niż 200 sesjach trend otrzymuje zero. Cały wynik techniczny jest niedostępny, jeśli brakuje sezonowości lub ostatnia świeca jest starsza niż 10 dni kalendarzowych.

% TODO(thesis): Add the detailed seasonality and SMA200 equations with a short numerical example.

\section{Sygnał z rejestru KNF}

Ranking porównuje najnowszą obserwację spółki z najnowszą obserwacją tego samego ISIN z wcześniejszego dnia UTC. Wzrost sumy ujawnionych pozycji daje $-1$, spadek $+1$, a niezmieniona suma zero. Dwie puste listy również dają zero; brak pary obserwacji oznacza \texttt{null}.

% TODO(thesis): Add one worked example using two archived KNF observations.

\section{Wynik końcowy}

Wynik końcowy ma postać
\begin{equation}
    Q=F+T+N,
\end{equation}
gdzie $F$ oznacza wynik fundamentalny, $T$ techniczny, a $N$ sygnał z KNF. Wagi wynoszą 1/1/1, bez normalizacji sumy. Brak dowolnej rodziny daje brak wyniku końcowego; liczby są sortowane malejąco, braki trafiają na koniec, a remisy są porządkowane według kodu spółki.

\cleardoublepage
\chapter{Implementacja}
\label{chap:implementation}

\section{Uruchomienie}

Instalacja z katalogu repozytorium i uruchomienie programu mają postać:
\begin{verbatim}
cargo install --path importer --locked
rankr
\end{verbatim}
Po odświeżeniu danych strona jest dostępna pod adresem \url{http://127.0.0.1:3000/}. Program nie wymaga argumentów ani konfiguracji bazy danych.

\section{Przetwarzanie danych}

Parsery sprawdzają metadane i układ odpowiedzi, a moduły obliczeniowe pracują na odczytanych strukturach. Niepoprawny plik jest zgłaszany i pomijany. Do rankingu wybierane są najnowsze dostępne obserwacje według czasu pobrania.

% TODO(thesis): Add selected implementation examples for parsing, incremental history updates and HTML rendering.

\cleardoublepage
\chapter{Testy i walidacja}
\label{chap:tests}

\section{Sprawdzenie programu}

Testy obejmują parsery, zapis, progi punktacji, łączenie rodzin i prezentację brakujących wartości. Korzystają z kontrolowanych danych oraz lokalnego serwera testowego. Zestaw można uruchomić poleceniem:
\begin{verbatim}
cargo test --workspace --locked --offline
\end{verbatim}

% TODO(thesis): Document the final test cases and a repeatable check of the complete local workflow.

\section{Ograniczenia}

Wagi rodzin są równe, ale ich zakresy są różne: fundamenty mają od $-5$ do $+5$ (w sektorze finansowym od $-1$ do $+1$), wynik techniczny od $-2$ do $+2$, a sygnał KNF od $-1$ do $+1$. Fundamenty mogą być sumą niepełnego zestawu wskaźników; nie obejmują jeszcze wyceny ani dynamiki rok do roku. Rejestr KNF pokazuje tylko ujawnione pozycje, więc pusta lista nie dowodzi braku krótkiej sprzedaży ani optymizmu inwestorów.

Program korzysta ze stałej mapy spółek i najnowszych zapisanych danych, bez odtwarzania historycznych składów indeksu i dostępności raportów. Odcięcie przyszłych świec nie wystarcza więc do wiarygodnego odtworzenia całego rankingu z przeszłości. Dopisywanie nowych cen nie odświeża też dawnych korekt dostawcy, a kompletność sesji nie jest sprawdzana względem niezależnego kalendarza giełdowego.

\section{Ocena przydatności rankingu}

Sprawdzenie kodu nie rozstrzyga, czy wyżej ocenione spółki osiągają później lepsze zwroty. Potrzebne jest zachowywanie rankingów przed okresem pomiaru i porównanie ich z późniejszymi wynikami spółek. Aplikacja nie zapisuje obecnie automatycznej historii rankingów.

% TODO(thesis): Define the observation period, return calculation, benchmark and treatment of tied scores.
% TODO(thesis): Add actual evaluation results with sample size and data coverage; do not infer 6-12 month performance from a shorter observation.

\cleardoublepage
\chapter{Podsumowanie}

Przygotowano lokalną aplikację, która pobiera dane finansowe, historię cen i ujawnione pozycje krótkie, oblicza trzy rodziny punktów oraz przedstawia ich sumę w przeglądarce. Obliczenia są oddzielone od pobierania danych i korzystają z lokalnych plików JSON. Dalszego opracowania wymaga ocena przydatności uzyskanej kolejności spółek.

% TODO(thesis): Complete the conclusions after the evaluation and final implementation scope are settled.

\cleardoublepage
\phantomsection
\addcontentsline{toc}{chapter}{Bibliografia}
\begin{thebibliography}{99}

\bibitem{stocks}
U.S. Securities and Exchange Commission, Investor.gov.
\newblock \textit{Stocks -- FAQs}.
\newblock \url{https://www.investor.gov/introduction-investing/investing-basics/investment-products/stocks}.
\newblock Dostęp: 28.10.2026.

\bibitem{gpwbenchmark}
GPW Benchmark.
\newblock \textit{Karta indeksu WIG20}.
\newblock \url{https://gpwbenchmark.pl/karta-indeksu?isin=PL9999999987}.
\newblock Dostęp: 28.10.2026.

\bibitem{gics}
MSCI.
\newblock \textit{The Global Industry Classification Standard (GICS)}.
\newblock \url{https://www.msci.com/indexes/index-resources/gics}.
\newblock Dostęp: 28.10.2026.

\bibitem{yahoo}
Yahoo Help.
\newblock \textit{What is the adjusted close?}
\newblock \url{https://in.help.yahoo.com/kb/adjusted-close-sln28256.html}.
\newblock Dostęp: 28.10.2026.

\bibitem{statements}
U.S. Securities and Exchange Commission.
\newblock \textit{Beginners' Guide to Financial Statement}.
\newblock \url{https://www.sec.gov/about/reports-publications/investorpubsbegfinstmtguide}.
\newblock Dostęp: 28.10.2026.

\bibitem{knfshorts}
Komisja Nadzoru Finansowego.
\newblock \textit{Rejestr krótkiej sprzedaży}.
\newblock \url{https://www.knf.gov.pl/podmioty/Rejestry_i_Ewidencje/rks?articleId=56674&p_id=18}.
\newblock Dostęp: 28.10.2026.

\bibitem{rust}
Rust Project Developers.
\newblock \textit{The Rust Programming Language: Packages and Crates}.
\newblock \url{https://doc.rust-lang.org/book/ch07-01-packages-and-crates.html}.
\newblock Dostęp: 28.10.2026.

\bibitem{axum}
Autorzy biblioteki Axum.
\newblock \textit{Dokumentacja Axum, wersja 0.8.9}.
\newblock \url{https://docs.rs/axum/0.8.9/axum/}.
\newblock Dostęp: 28.10.2026.

\bibitem{tokio}
Autorzy Tokio.
\newblock \textit{Tokio Tutorial}.
\newblock \url{https://tokio.rs/tokio/tutorial}.
\newblock Dostęp: 28.10.2026.

\bibitem{serdejson}
Autorzy biblioteki serde\_json.
\newblock \textit{Dokumentacja Serde JSON, wersja 1.0.151}.
\newblock \url{https://docs.rs/serde_json/1.0.151/serde_json/}.
\newblock Dostęp: 28.10.2026.

\bibitem{sqlite}
SQLite.
\newblock \textit{About SQLite}.
\newblock \url{https://www.sqlite.org/about.html}.
\newblock Dostęp: 28.10.2026.

\bibitem{decimal}
Autorzy biblioteki rust\_decimal.
\newblock \textit{Dokumentacja rust\_decimal, wersja 1.43.0}.
\newblock \url{https://docs.rs/rust_decimal/1.43.0/rust_decimal/}.
\newblock Dostęp: 28.10.2026.

\end{thebibliography}

\cleardoublepage
\begingroup
\small
\linespread{1.0}\selectfont
\thispagestyle{empty}
\setlength{\parindent}{0in}
\setlength{\parskip}{0em}
\phantomsection
\addcontentsline{toc}{chapter}{Streszczenie}

\textbf{POLITECHNIKA RZESZOWSKA} \hfill Rzeszów \ThesisYear \\
\textbf{im. Ignacego Łukasiewicza} \\
\textbf{WYDZIAŁ MATEMATYKI I FIZYKI STOSOWANEJ} \\

\begin{center}
    \textbf{STRESZCZENIE PRACY DYPLOMOWEJ}
\end{center}

\textbf{Tytuł:} \ThesisTitle \\
\textbf{Autor:} \AuthorName \\
\textbf{Promotor:} \SupervisorName \\
\textbf{Słowa kluczowe:} GPW, WIG20, scoring, ranking, Rust, dane finansowe
\\

Praca opisuje lokalną aplikację webową \textit{rankr} w języku Rust do punktowego porównywania spółek z mapy WIG20 projektu. Program pobiera dane finansowe GPW / Notoria, dzienne ceny Yahoo Finance i ujawnione pozycje krótkie KNF. Zapisuje dane lokalnie, oblicza trzy rodziny punktów i przedstawia ich sumę w tabeli HTML obsługiwanej przez Axum. Opis obejmuje źródła, architekturę, reguły punktacji i ograniczenia. Ocena związku rankingu z późniejszymi zwrotami pozostaje do wykonania.

\vfill
\noindent\rule{\textwidth}{2pt}

\textbf{RZESZOW UNIVERSITY OF TECHNOLOGY} \hfill Rzeszów \ThesisYear \\
\textbf{FACULTY OF MATHEMATICS AND APPLIED PHYSICS} \\

\begin{center}
    \textbf{DIPLOMA THESIS ABSTRACT}
\end{center}

\textbf{Title:} \ThesisTitleEnglish \\
\textbf{Author:} \AuthorName \\
\textbf{Supervisor:} Mariusz Startek, PhD \\
\textbf{Key words:} Warsaw Stock Exchange, WIG20, scoring, ranking, Rust, financial data
\\

This thesis describes \textit{rankr}, a local web application written in Rust for scoring companies listed in the project's WIG20 map. The program collects GPW / Notoria financial data, Yahoo Finance daily prices and disclosed KNF net short positions. It stores data locally, calculates three score families and displays their sum in an HTML table served by Axum. The thesis outlines data sources, architecture, scoring rules and limitations. The relationship between the ranking and subsequent returns remains to be evaluated.

\vfill
\clearpage
\endgroup

\cleardoublepage
\thispagestyle{empty}
\phantomsection
\addcontentsline{toc}{chapter}{Wykaz wykorzystania narzędzi GenAI}

\begin{center}
    Wykaz obszarów i narzędzi GenAI \\
    wykorzystanych przez autora w trakcie wykonywania pracy dyplomowej
\end{center}

\begin{table}[htbp]
    \centering
    \begin{tabularx}{\linewidth}{|c|X|X|}\hline
       \textbf{Lp.} & \multicolumn{1}{c|}{\textbf{Obszar wykorzystywania}} & \multicolumn{1}{c|}{\textbf{Narzędzie}} \\ \hline
       1 & Przygotowanie i redakcja fragmentów pracy, korekta techniczna oraz porównanie opisu z kodem projektu. & ChatGPT / Codex \\ \hline
       2 & Wsparcie w planowaniu zakresu, implementacji pobierania danych, punktacji i strony rankingu oraz przygotowaniu testów. & ChatGPT / Codex \\ \hline
       3 & Wsparcie w analizie dokumentacji technicznej i przygotowaniu prezentacji postępu prac. & ChatGPT / Codex \\ \hline
    \end{tabularx}
\end{table}

\end{document}
